%\thispagestyle{empty}

Compreender o padrão de radiação das antenas é essencial para otimizar os sistemas de comunicação sem fios. Esse parâmetro descreve como a antena irradia ou recebe energia eletromagnética no espaço, tais como a área de cobertura, a intensidade do sinal e as interferências. Conhecer o padrão de radiação permite posicionar e orientar melhor as antenas, minimiza as perdas de energia e melhora o desempenho geral do sistema.\\
Embora os softwares de simulação eletromagnética (EMS) possam estimar o padrão de radiação durante a fase de projeto, os resultados práticos muitas vezes diferem das simulações devido a influências ambientais. Reflexões, difrações e absorções podem distorcer significativamente o que é idealizado. Como resultado, compreender o padrão de radiação real é crucial para garantir um desempenho confiável em aplicações práticas.\\
No entanto, os equipamentos de medição de antenas comerciais são normalmente muito caros, o que os torna impraticáveis para laboratórios educacionais. Neste contexto, este trabalho propõe o desenvolvimento de uma configuração experimental de baixo custo para medir padrões de radiação de antenas na banda Industrial, Científica e Médica (ISM). O sistema permite o posicionamento automático de uma antena de prova segundo dois eixos, habiltando a medição de padrão coplanar e cruzada assim como azimute e/ou elevação. O controlo dos dois graus de liberdade, assim como a aquisição e visulização dos dados, é realizado por um programa desenhado para correr no software MATLAB. A validação experimental, com uma antena Yagi-Uda, mostrou boa concordância entre os resultados medidos e simulados. Os resultados validam a operação do sistema desenvolvido.



\noindent {\bf Palavras-chave:} Padrão de Radiação de Antenas, Radiação Eletromagnética, Mecatrónica, Medição de Baixo Custo, Desempenho de Antenas.


%\vfill
%\pagebreak
%\mbox{}
%\vfill
%\pagebreak